% ID: SEL-D0014-BC-E05
% ESTADO: validada. S015; tres pasadas del mismo autor, sin segundo revisor.
\textbf{Modelo sin reemplazamiento.}\par\nopagebreak[4]
Se eligen dos llaves distintas de las $90$. Tras la primera elección quedan $89$, y la composición depende de si la primera funciona. No son selecciones independientes.

\textbf{a) Ambas funcionan.}\par\nopagebreak[4]
\[
\boxed{P(\text{dos buenas})=\frac{60}{90}\frac{59}{89}
=\frac{118}{267}\approx0.441948.}
\]
\textbf{b) Al menos una funciona.}\par\nopagebreak[4]
Es el complemento de escoger las dos que no funcionan:
\[
\boxed{P(\text{alguna buena})=1-\frac{30}{90}\frac{29}{89}
=\frac{238}{267}\approx0.891386.}
\]
``Alguna'' incluye el caso de que funcionen las dos.

\textbf{c) Exactamente una funciona.}\par\nopagebreak[4]
Hay dos órdenes disjuntos: buena/mala y mala/buena:
\[
\boxed{P(\text{una buena})=\frac{60}{90}\frac{30}{89}
+\frac{30}{90}\frac{60}{89}=\frac{40}{89}\approx0.449438.}
\]
Como comprobación, también se obtiene $\binom{60}{1}\binom{30}{1}/\binom{90}{2}$. La probabilidad de al menos una es la suma de los resultados a y c; junto a ninguna, las tres posibilidades suman uno.
